\documentclass[aps,reprint]{revtex4-1}
% Engine-specific settings
% Detect pdftex/xetex/luatex, and load appropriate font packages.
% This is inspired by the approach in the iftex package.
% pdftex:
\ifx\pdfmatch\undefined
\else
    \usepackage[T1]{fontenc}
    \usepackage[utf8]{inputenc}
\fi
% xetex:
\ifx\XeTeXinterchartoks\undefined
\else
    \usepackage{fontspec}
    \defaultfontfeatures{Ligatures=TeX}
\fi
% luatex:
\ifx\directlua\undefined
\else
    \usepackage{fontspec}
\fi
% End engine-specific settings
\usepackage[english]{babel}
\usepackage{csquotes}
% \usepackage[backend=biber, sortcites]{biblatex}
\usepackage{url}
\usepackage{textcomp}
\usepackage[usenames,dvipsnames,svgnames, table]{xcolor}
\usepackage[font={scriptsize}]{caption}
\usepackage{amsmath} \usepackage{amsthm} \usepackage{amsfonts}
\usepackage{amssymb}
\usepackage{enumerate}
\usepackage{tikz} \usepackage{float}
\usepackage[procnames]{listings}
\usepackage{pstool} \usepackage{pgfplots}
\usepackage{wrapfig} \usepackage{graphicx} \usepackage{epstopdf}
\usepackage{afterpage}
\usepackage{physics}
\usepackage{multirow}
\usepackage{gensymb}
\usepackage{algorithm}
\usepackage{microtype}
\usepackage[noend]{algpseudocode}
\usepackage{xcolor,colortbl}
\usepackage{microtype}
\usepackage{geometry}
\usepackage{hyperref}
\usepackage{graphicx}
\usepackage{caption}
\usepackage{subcaption}
\usepackage{lipsum}
% \usepackage{pythontex}
% \usepackage{authblk}
\usepackage{nth}
\usepackage{siunitx}
% \usepackage[toc,page]{appendix}
\floatstyle{plaintop}
\restylefloat{table}

% SET LISTING STUFF (FRED)
\definecolor{dkgreen}{rgb}{0,0.6,0}
\definecolor{dred}{rgb}{0.545,0,0}
\definecolor{dblue}{rgb}{0,0,0.545}
\definecolor{lgrey}{rgb}{0.9,0.9,0.9}
\definecolor{gray}{rgb}{0.4,0.4,0.4}
\definecolor{darkblue}{rgb}{0.0,0.0,0.6}
\lstdefinelanguage{cpp}{
      backgroundcolor=\color{lgrey},
      basicstyle=\footnotesize \ttfamily \color{black} \bfseries,
      breakatwhitespace=false,
      breaklines=true,
      captionpos=b,
      commentstyle=\color{dkgreen},
      deletekeywords={...},
      escapeinside={\%*}{*)},
      frame=off,
      language=C++,
      keywordstyle=\color{purple},
      morekeywords={BRIEFDescriptorConfig,string,TiXmlNode,DetectorDescriptorConfigContainer,istringstream,cerr,exit},
      identifierstyle=\color{black},
      stringstyle=\color{blue},
      % numbers=left,
      % numbersep=5pt,
      numberstyle=\tiny\color{black},
      rulecolor=\color{black},
      showspaces=false,
      showstringspaces=false,
      showtabs=false,
      stepnumber=1,
      tabsize=5,
      title=\lstname,
    }

% Custom commands
\newcommand{\unit}[1]{\:\mathrm{#1}}
\newcommand{\noref}[1]{\hyperref[#1]{\ref*{#1}}}
\newcommand{\nonref}[1]{\hyperref[]{\ref*{#1}}}
\newcommand\blankpage{%
  \null
  \thispagestyle{empty}%
  \addtocounter{page}{-1}%
  \newpage}

% Default fixed font does not support bold face
\DeclareFixedFont{\ttb}{T1}{txtt}{bx}{n}{7} % for bold
\DeclareFixedFont{\ttm}{T1}{txtt}{m}{n}{7}  % for normal

\newcommand\numberthis{\addtocounter{equation}{1}\tag{\theequation}}
\DeclareCaptionFont{white}{\color{white}}
\DeclareCaptionFormat{listing}{\colorbox{gray}{\parbox{\columnwidth}{#1#2#3}}}
\captionsetup[lstlisting]{format=listing,labelfont=white,textfont=white}


% Biber for references
% \bibliographystyle{aipauth4-1}

\begin{document}
\sisetup{detect-all}
\title{Eigenvalue stuff}
\author{Erlend Lima}
\author{Frederik J. Mellbye}
\author{Aram Salihi}
\author{Halvard Sutterud}
\affiliation{University of Oslo, Oslo, Norway \\ Source code available at: \url{https://github.com/Caronthir/FYS3150/tree/master/Project2}}
\date{\today}

\begin{abstract}
  This is really abstract
\end{abstract}
\maketitle
\tableofcontents
\makeatletter
\let\toc@pre\relax
\let\toc@post\relax
\makeatother 

\newpage

% % Get all of the data from the script into a dictionary
% \begin{pycode}
% from subprocess import Popen, PIPE
% from latextools import untag_all
% p = Popen(['python', '../analysis/analyze.py', '../cpp'],
% stdout=PIPE, stderr=PIPE)
% output, err = p.communicate()
% data = output.decode('utf8')
% latex_data = untag_all(data)
% \end{pycode}
%
% \begin{table}[ht]
%   \centering
%   \pyc{print(latex_data['error_table'])}
%   \caption{Summary of the errors}
%   \label{tab:error}
% \end{table}
%
% \begin{figure}[ht]
%   \centering
%   \includegraphics[scale=0.5]{figures/function.eps}
%   \caption{\label{fig:general} Caption}
% \end{figure}
\section{Introduction}
\label{sec:introduction}
\section{Theory}
\label{sec:theory}
\subsection{Idiotisk navn for å understreke at dette er midlertidig}
A unitary transformation preserves orthogonality and the dot product. To see
this, consider an orthogonal basis \(\{\vb{v}_{i}\}\) and an unitary
transformation \(U\). Defining \(\vb{w}_{i}=U\vb{v}_{i}\), the dot product can
be computed as
\begin{align*}
  \vb{w}_{i}\vdot \vb{w}_{j} &= \vb{w}_{i}\vb{w}_{j}^{T}\\
                             &= \left( U\vb{v}_{i} \right)\left( U\vb{v}_{j} \right)^{T}\\
                             &= U\vb{v}_{i}\vb{v}_{j}^{T}U^{T}\\
                             &= U\delta_{ij}U^{T}\\
                             &= \delta_{ij}UU^{T}\\
                             &= \delta_{ij}
\end{align*}
where the last step follows since \(U\) is unitary. 
\section{Method}
\label{sec:method}
\section{Results}
\label{sec:results} 
\section{Discussion}
\label{sec:discussion}
\section{Conclusion}
\label{sec:conclusion}
\bibliography{references}
\blankpage
\appendix
\blankpage
\end{document}
